\section{Setting}
\label{sec:setting}

\subsection{Critical rate and gap}

Let $\dot x = f(x)$ be a Hamiltonian flow and let the observer see $Y = X + \sigma Z$, where $X$ has density $\rho$ and $Z$ is standard normal. The coarse-grained density is $\varrho = \rho * N(0, \sigma^2 I)$. We write
\begin{equation}
s = \nabla\log\varrho, \qquad H = -\nabla^2\log\varrho, \qquad m = \E[X \mid Y = y],
\end{equation}
and $\E$ for the expectation over $Y \sim \varrho$. Tweedie's formula gives $m = y + \sigma^2 s$ and $\Cov(X \mid y) = \sigma^2 I - \sigma^4 H$~\cite{efron2011}. When the width grows as $\sigma(t) = \sigma_0 e^{rt}$, the coarse-grained entropy obeys Eq.~\eqref{eq:SC} with the critical rate~\cite{namba2026arrow,namba2026learned}
\begin{equation}
R = \frac{\E\big[s \cdot \E[f(X) \mid Y]\big]}{\sigma^2\,\E|s|^2}.
\label{eq:R}
\end{equation}
Let $\kappa$ be the largest eigenvalue of the symmetric part of $Df$ over phase space. We study the \emph{gap in units of the resolution},
\begin{equation}
g = \frac{\kappa - R}{\sigma}.
\end{equation}
For every Gaussian state $g \geq 0$, and for the best Gaussian state of the pendulum $g \to 1$ as $\sigma \to 0$~\cite{namba2026gap}.

\subsection{The top of the pendulum}

For $H = p^2/2 - \cos q$ we use $z = q - \pi$, so that near the unstable top the flow is
\begin{equation}
\dot z = p, \qquad \dot p = \sin z,
\end{equation}
with $\kappa = 1$. The unstable and stable directions of the linearised flow are $V = (1,1)/\sqrt2$ and $W = (1,-1)/\sqrt2$ in the $(z,p)$ plane, with coordinates $v = (z + p)/\sqrt2$ and $w = (z - p)/\sqrt2$. Subscripts $v, w, z, p$ denote components in these directions, and
\begin{gather}
J_{vv} = \E[s_v^2], \qquad J_{ww} = \E[s_w^2], \nonumber\\
a = \sigma^2 J_{vv}, \qquad d = \sigma^2(J_{vv} + J_{ww}).
\end{gather}
We call $a$ the \emph{sharpness} of the state: it is the Fisher information along the unstable direction in units of $\sigma^{-2}$. Since $\varrho$ is a $\sigma$-blur, $0 < a \leq 1$.

The pendulum lives on a cylinder, and the analysis near the top needs the state to stay within about one fundamental domain. We use two moment conditions,
\begin{gather}
\text{(H1)} \qquad \E_Q[y_z^2] \leq \pi^2 + \sigma^2/2, \qquad Q = \varrho^2/\|\varrho\|^2, \nonumber\\
\text{(H2)} \qquad \Var(X_z) \leq \pi^2 .
\label{eq:H0}
\end{gather}
Both follow from $\rho = 0$ outside $|z| \leq \pi$ (for (H1) see Appendix~\ref{app:split}), but neither requires it: a Gaussian state, whose density never vanishes, satisfies both when it is narrow compared with $\pi$. Each result below states which condition it uses. Under (H2) the Cram\'er--Rao inequality $J \succeq \Cov(Y)^{-1}$~\cite{cover2006}, together with $e_w\cdot e_z = 1/\sqrt2$, gives
\begin{equation}
J_{ww} \geq \frac{1}{2\Var(Y_z)} \geq \frac{1}{2(\pi^2 + \sigma^2)},
\label{eq:CR}
\end{equation}
so that $\ln\big(1/(\sigma^2 J_{ww})\big) \leq 2\ln(1/\sigma) + \ln\big(2(\pi^2 + \sigma^2)\big)$: the information along the stable direction cannot fall below a floor set by the confinement.

\subsection{A formula without cancellation}

In Eq.~\eqref{eq:R} the numerator contains two terms of order one that cancel up to the gap, which is of order $\sigma$, or much smaller in the states of Sec.~\ref{sec:strips}. The following form removes the cancellation analytically.

\begin{proposition}[Gap formula]
\label{prop:formula}
For the pendulum near the top,
\begin{equation}
\sigma g = \frac{2\sigma^2 J_{ww} + N}{d}, \qquad N = \E[s_p D] = -\E[\partial_p D],
\label{eq:formula}
\end{equation}
where $D(y) = \E[z - \sin z \mid Y = y]$ is the nonlinear part of the conditional velocity. Equivalently,
\begin{equation}
N = -\sigma^{-2}\,\E\big[\Cov(X_p,\, X_z - \sin X_z \mid Y)\big].
\label{eq:Ncov}
\end{equation}
\end{proposition}

\begin{proof}
$\E[f(X)\mid y] = (m_p,\, m_z - D)$ in the $(z,p)$ components, so $s\cdot\E[f\mid y] = s_z m_p + s_p m_z - s_p D$. With $m = y + \sigma^2 s$, and $\E[s_z y_p] = \E[s_p y_z] = 0$ by integration by parts, one finds $\E[s_z m_p + s_p m_z] = 2\sigma^2\E[s_z s_p] = \sigma^2(J_{vv} - J_{ww})$. Inserting this in Eq.~\eqref{eq:R} with $\kappa = 1$ and $\sigma^2\E|s|^2 = d$ gives Eq.~\eqref{eq:formula}. The second form of $N$ follows from $\partial_y \E[h(X)\mid y] = \sigma^{-2}\Cov(h(X), X \mid y)$.
\end{proof}

Equation~\eqref{eq:formula} has a simple reading. The gap is the Fisher information along the stable direction, $2\sigma^2 J_{ww}$, plus the effect $N$ of the nonlinearity $z - \sin z$ of the pendulum, both divided by the total information $d$. For a sharp state, $d \approx a$. The gap is small when the state carries little information along the stable direction \emph{and} feels little nonlinearity, while being sharp along the unstable one. The rest of this article is about how far these demands can be met at the same time.

In the numerical searches we also use the focus ratio $\Upsilon = \sigma J_{ww}/j$, with $j = \E[\min(\sigma^2 (H_{vv})_+,\, 1)]$, and restrict to $\Upsilon \leq 1/2$. This keeps the searches among focused states, in which the information along the stable direction is small compared with the sharpness $a$ and $a/d$ is close to one. All states discussed below satisfy it; none of the proofs requires it. We also write $t = \sigma^2$ for the variance of the blur when the blur is varied, as in Sec.~\ref{sec:removing}.
